\begin{table}[htbp]
\centering
\small
\caption{Sample construction.}
\label{tab:t2_sample}
\begin{tabular}{lllrllr}
\toprule
City & Manufacturing MSMEs & Artefacts removed & % removed & Analysis sample & Unique locations & Enterprises per location \\
\midrule
Patna & 14,316 & 1,509 & 10.500000 & 12,807 & 6,312 & 2.030000 \\
Indore & 67,270 & 18,386 & 27.300000 & 48,884 & 19,854 & 2.460000 \\
Varanasi & 10,973 & 1,684 & 15.300000 & 9,289 & 3,923 & 2.370000 \\
Kochi & 43,464 & 15,705 & 36.100000 & 27,759 & 12,222 & 2.270000 \\
Kolkata & 26,129 & 1,277 & 4.900000 & 24,852 & 13,126 & 1.890000 \\
Total & 162,152 & 38,561 & 23.800000 & 123,591 & 55,437 & 2.230000 \\
\bottomrule
\end{tabular}
\begin{tablenotes}\footnotesize
\item Artefacts identified by two independent rules (Section 3.5): at least 20 enterprises at one coordinate spanning at least five postal codes, or at least eight ISIC divisions. Registry co-location means effective spatial sample size is roughly half the enterprise count.
\end{tablenotes}
\end{table}
